\documentclass[a4paper, 10pt]{article}

\usepackage[margin = 1in]{geometry} % for spacing around
\usepackage{graphicx} % for including images in your pdfs
\usepackage{xcolor} % for including colors in your pdf
\usepackage{soul} % for text decoration
\usepackage[utf8]{inputenc} % for encoded text
\usepackage[T1]{fontenc}
\usepackage{setspace} % for setting different line spacings between paragrafs.
\usepackage{enumerate} % for letting us get more detailed enumerate lists
\usepackage{multirow} % to let us combine more rows together
\usepackage{colortbl} % for decorating tables
\usepackage{amsmath} % used for representing more complicated math displays
\usepackage{supertabular}
\usepackage{longtable} % both of these packages are used to making really big tables
\usepackage{wrapfig} % allows us to wrap text around figures
\usepackage{fancyhdr} % for making fancy headers
%\usepackage{bibtex} % for making better bibliographies
\usepackage[pdftex]{hyperref} % for letting us make links
\usepackage{lscape} % Allows us to flip from portrait to landspace
\usepackage{tikz} % for high detailed drawing
\usepackage{multicol} % To put things side by side
\usepackage{rotating} % For rotating objects
% \usepackage{draftwatermark} % For adding watermarks
\usepackage{MnSymbol} % for using multiple symbols
\usepackage{mathtools} % Used for more math symbols
\usepackage{xfrac} % For more complciated fractions and to add derivitives
\usepackage{enumitem} % for better enum lists
\usepackage{tcolorbox} % for adding colored text boxes
\usepackage{bm} % Adding bold text to math inputs
\usepackage{pgfplots} % Used for plotting functions
\usepackage{wallpaper} % For adding wallpapers
\usepackage{listings} % Required for insertion of code
\usepackage{eso-pic} % Required for placing background pictures
\usepackage{ifthen} % Required for conditional statements

% Define the Haskell language style for listings
\lstdefinelanguage{Haskell}{
    morekeywords={class, data, deriving, do, else, if, import, in, infixl, infixr, instance, let, module, newtype, of, then, type, where, _, ::, =, \, |, <-, ->, @, ~, =>},
    sensitive=true,
    morecomment=[l]{--},
    morecomment=[s]{\{-}{-\}},
    morestring=[b]",
    stringstyle=\color{red},
    identifierstyle=\color{blue},
    keywordstyle=\color{green},
    commentstyle=\color{gray},
    basicstyle=\ttfamily, % Use monospaced font for code
    showstringspaces=false, % Don't show spaces in strings as special underscores
    numbers=left, % Line numbers on the left
    numberstyle=\tiny\color{gray}, % Line numbers are grey and small
    stepnumber=1 % Line numbers go in steps of 1
}

% Setting up the default image path
\graphicspath{{./Images/}}

% Implementing authro details
\title{Introduction to Haskell Programming Language}
\author{Emre Arapcic-Uevak}
\date{}

% Setting up the fancy page style
\fancypagestyle{customStyle}{
	\lhead{} \chead{} \rhead{}
	\lfoot{} \cfoot{\thepage} \rfoot{}
	\renewcommand{\headrulewidth}{0pt}
	\renewcommand{\footrulewidth}{1pt}
}
\pagestyle{customStyle}

% Setting up hyperref options
\hypersetup {
	colorlinks = false,
	citecolor = black,
	filecolor = blue,
	linkcolor = blue,
	urlcolor = blue,
	pdftex
}

% Custom commands


\begin{document}
	\maketitle
    \lstset{language=Haskell}
	\vspace{5mm}
	
	\begin{abstract}
		\begin{center}
			This presentation introduces the Haskell programming language, covering its history, key features, data structures, comparisons with other languages, and providing simple code examples. The audience will gain insights into the fundamentals of Haskell and its applications in modern programming.
		\end{center}
	\end{abstract}

	\ThisCenterWallPaper{.4}{Images/IUS Logo.png}

	\newpage % Start new page after the first page
    \AddToShipoutPictureBG{%
        \ifthenelse{\value{page}=1}{}{%
            \begin{tikzpicture}[remember picture,overlay]
                \node [opacity=0.3, at=(current page.center)] {
                    \includegraphics[width=0.5\paperwidth]{Haskell-Logo}
                };
            \end{tikzpicture}
        }
    }

	\tableofcontents
	\pagebreak

	\section{Introduction to Haskell}
        \noindent Haskell is a purely functional programming language, distinguished by its strong static typing, high level of abstraction, and lazy evaluation. Inspired by the principles of lambda calculus, Haskell emphasizes functions without side effects and immutable data. Named after the logician Haskell Curry, it stands apart from imperative programming languages by treating computation as the evaluation of mathematical functions, thereby avoiding changing-state and mutable data.
	\section{History and Evolution of Haskell}
        \subsection{The Foundational Principles}
            Haskell's journey begins with its roots deeply embedded in lambda calculus, a sophisticated mathematical system devised by Alonzo Church in 1932. This system, focusing on function abstraction and application, forms the core of Haskell's functional programming ethos.

        \subsection{The Dawn of Functional Programming}
            Tracing its lineage to the 1960s, Haskell shares its heritage with LISP, the first functional language, born in 1959. However, Haskell emerged as a response to the limitations of its predecessors like ML, Hope, and Miranda, which, despite their innovations, struggled in practical application and widespread adoption.

        \subsection{The Formation of Haskell}
            Haskell's story officially starts at the FPCA '87 conference in Portland, Oregon. Here, a visionary group of academics united to forge a new path in programming language design. Their collaboration led to the first Haskell report in 1990, laying out the motivations, aspirations, and the unique nature of Haskell.

        \subsection{Milestones in Haskell’s Journey}
            Haskell's evolution saw pivotal moments, such as the release of the "Gentle Introduction to Haskell" tutorial and the advent of its essential tools – the GHC compiler and GHCI interpreter in 1992. At a seminal meeting at Yale in 1988, Haskell's goals were crystallized, focusing on research, teaching, and the development of robust systems.

        \subsection{A Tribute to Haskell B. Curry}
            Named to honor the eminent logician and mathematician Haskell B. Curry, the language reflects his groundbreaking work. Successive versions, including the Haskell Report 1.1 and version 1.2, introduced nuanced enhancements, shaping Haskell into the language we recognize today.

        \subsection{Haskell in the Digital Age}
            Marking its online presence, Haskell took to the internet in 1994 with haskell.org, a domain that continues to be a hub for enthusiasts. The release of Haskell version 1.3 in 1996 was a landmark, introducing significant refinements and capabilities.

        \subsection{Haskell 98 - A Commitment to Stability}
            The publication of "The Haskell 98 Report: Language and Libraries" in 1999 was a declaration of Haskell's maturity and stability, broadening its appeal and application. This version, later available freely online, cemented Haskell's place in the programming world.

        \subsection{The 21st Century and Haskell Prime}
            With the turn of the century, Haskell faced the need for evolution. The emergence of Haskell Prime (Haskell$'$) signified a modernized, incrementally developed version, with Haskell 2010 being its first notable iteration.

        \subsection{Haskell Today}
            Haskell's narrative continues to unfold, with the Haskell Foundation, established in 2020, championing its development. The support from numerous companies ensures Haskell's enduring presence and evolution in the realm of functional programming.

    \pagebreak
	\section{Features of the Language}
        Haskell is distinguished by a range of unique and powerful features, each contributing to its robustness as a functional programming language.

        \subsection{Purely Functional Programming}
            Haskell is a purely functional language, ensuring functions are free from side effects. This approach simplifies both debugging and testing, as the same input always yields the same output.

        \subsection{Strong, Static Type System with Type Inference}
            The language boasts a robust and static type system, with type checking at compile time to minimize runtime errors. Type inference in Haskell allows for less verbose code while maintaining type safety.

        \subsection{Lazy Evaluation}
            Haskell employs lazy evaluation, delaying computations until their results are needed. This leads to efficient memory utilization and the ability to define infinite data structures.

        \subsection{Immutability}
            Variables in Haskell are immutable. Once a value is assigned to a variable, it cannot be altered, which aids in reducing side effects and simplifying code maintenance.

        \subsection{First-Class Functions}
            Functions in Haskell are treated as first-class citizens, capable of being passed as arguments, returned from other functions, and assigned to variables.

        \subsection{Pattern Matching}
            Haskell offers advanced pattern matching capabilities, allowing for concise and readable code through the direct decomposition of data structures.

        \subsection{High-Level Abstractions}
            The language supports high-level abstractions such as monads, facilitating the management of complex computational patterns like side effects.

        \subsection{Concurrent and Parallel Programming}
            Haskell excels in concurrent and parallel programming, streamlining the development of efficient multi-core processor applications.

        \subsection{Extensive Standard Library}
            Haskell is equipped with a comprehensive standard library, covering a broad spectrum of functionalities, from basic data manipulation to sophisticated algorithms.

        \pagebreak

	\section{Structures of the Language}
        Haskell's design and architecture are characterized by several key structural elements, which make it unique and powerful among programming languages.

        \subsection{Syntax and Notation}
            Haskell features a clean and concise syntax that emphasizes readability and maintainability. It uses significant whitespace and layout rules, similar to Python, to define code structure, reducing the reliance on braces and semicolons.
            \textbf{Key Aspects of Haskell's Syntax:}
            \begin{itemize}
                \item \textit{Function Definition and Application:} Functions are defined using a straightforward syntax, where the function name is followed by its parameters, and the body of the function follows an equals sign. Function application is equally straightforward, requiring only the function name followed by its arguments, without parentheses or commas.
                \item \textit{Indentation-Sensitive:} Haskell's layout rule dictates that code grouped together (like the body of a function or a control structure) must be indented under the same column. This leads to a visually structured and clean code layout.
                \item \textit{Infix and Prefix Notation:} Haskell allows both infix and prefix notation for functions, enhancing flexibility and readability in different contexts. For example, operators like `+` are used infix, while most other functions are used in prefix form.
            \end{itemize}

            \textbf{Example of Haskell Syntax:}
            Consider a simple function definition in Haskell:

            \begin{lstlisting}
addNumbers :: Int -> Int -> Int
addNumbers x y = x + y
            \end{lstlisting}

            This function `addNumbers` takes two integers (`Int`) as inputs and returns their sum. The type declaration (`::`) is optional but helps in understanding the function's contract. Notice the absence of curly braces and semicolons, and how the function's body is simply defined after the equals sign.

            Haskell's syntax, while different from many mainstream languages, offers a level of expressiveness and clarity that is appreciated in complex software development.

        \subsection{Basic Types}

            Haskell's standard library includes several built-in basic types that are fundamental to the language. The most commonly used are:

            \begin{itemize}
                \item \texttt{Bool} for boolean values.
                \item \texttt{Char} for single Unicode characters.
                \item \texttt{String}, which is syntactically equivalent to a list of characters (\texttt{[Char]}).
                \item \texttt{Int} for fixed-precision integers, typically 64 bits in size on modern systems.
                \item \texttt{Integer} for arbitrary-precision integers, limited only by available memory.
                \item \texttt{Float} for single-precision floating-point numbers.
                \item \texttt{Double} for double-precision floating-point numbers.
            \end{itemize}

            It's worth noting that numeric literals in Haskell are polymorphic. They can represent any numeric type based on the context in which they are used.

            Here are some illustrative examples:

            \begin{lstlisting}
x :: Int
x = 42  -- Represents a fixed-size integer

y :: Double
y = 3.14  -- Represents a double-precision floating-point number

isTrue :: Bool
isTrue = true  -- Represents a boolean value

charExample :: Char
charExample = 'A'  -- Represents a single character

bigInt :: Integer
bigInt = 123456789012345678901234567890  -- Represents an arbitrary-size integer

someString :: String
someString = "This is a string"  -- Represents a string, equivalent to a list of characters
            \end{lstlisting}


        \subsection{Explicit and Implicit Declaration (Type Inferencing)}
            Haskell allows for both explicit and implicit type declarations. Explicit declaration is specifying the type of a variable or function using the double colon operator \texttt{::}, while implicit declaration leverages the compiler's type inference system.

            \begin{lstlisting}
myVar :: String
myVar = "CS305"
            \end{lstlisting}

            Implicit type declaration does not require specifying types, as the compiler deduces the type from the values used.

            \begin{lstlisting}
myVar = "CS305"  -- The compiler infers myVar as a String
            \end{lstlisting}

        \subsection{Comments}
            Comments are integral to writing clear and maintainable code. In Haskell, as in many programming languages, comments come in two forms: single-line and multi-line. Single-line comments are denoted by \texttt{--} and extend to the end of the line. Multi-line comments are enclosed within \texttt{\{-} and \texttt{-\}}, and uniquely, they can be nested within one another, which is a feature not commonly found in all programming languages.

            \lstset{language=Haskell}
            \begin{lstlisting}
-- This is a single-line comment

{-
This
is
a
multi-line
comment
which can be {- nested -} within another comment
-}
            \end{lstlisting}

            These commenting conventions facilitate the documentation of code and the temporary disabling of code blocks during debugging and development.

        \subsection{Immutability}

            Immutability is a cornerstone of Haskell's design, mandating that variables, once assigned, cannot be altered. This principle underpins the functional programming paradigm, where data stability and predictability are paramount.

            \lstset{language=Haskell}
            \begin{lstlisting}
someVar :: String
someVar = "Modify me"

-- Reassignment is not allowed and results in a compile-time error:
-- someVar = "modified"  -- This causes a multiple declaration error
            \end{lstlisting}

            Attempting to reassign `someVar` to a new value here would result in a compile-time error, specifically a "multiple declaration error," indicating that `someVar` has been assigned a value more than once.

            This approach contrasts with mutable languages, where variables can change state throughout their lifecycle. While immutability can lead to a more functional and safer codebase, it may require programmers to think differently, particularly when it comes to managing state and side effects. Some of the benefits include easier reasoning about code behavior and enhanced thread safety. However, it may also introduce overhead due to the creation of new instances and could necessitate a deeper understanding of recursion and state management techniques within Haskell's functional context.


    \subsection{Lists}
        A list in Haskell is a collection of elements of the same type, enclosed in square brackets and separated by commas. Haskell lists are zero-indexed and can be manipulated using a variety of functions, such as \texttt{head} for the first element, \texttt{tail} for all but the first element, and \texttt{length} for the number of elements.

        \lstset{language=Haskell}
        \begin{lstlisting}
nums :: [Int] -- This type declaration is optional
nums = [1, 2, 3, 4, 5]
        \end{lstlisting}

        Haskell's lazy evaluation model allows for lists to be potentially infinite. This feature enables defining lists without predetermined bounds, such as the list of all positive integers:

        \begin{lstlisting}
positiveIntegers = [1..] -- An infinite list of positive integers
        \end{lstlisting}

        Lazy evaluation ensures that operations on such infinite lists are still manageable, as elements are only computed as needed. For instance, list comprehensions can be used to create lists from existing lists:

        \begin{lstlisting}
squares = [x * x | x <- [1..10]] -- List of squares from 1 to 10
        \end{lstlisting}

        This list comprehension takes numbers 1 through 10 and returns a new list of their squares. Operations like these demonstrate the power and flexibility of Haskell's list handling capabilities.



        \subsection{Higher-Order Functions}
            In Haskell, it is possible to define functions that accept other functions as arguments, which classifies them as higher-order functions. A classic example of this is the \texttt{map} function, which applies a given function to each item in a list.
            Here's a simple example:

            \begin{lstlisting}
squareNum :: Int -> Int
squareNum n = n * n

nums :: [Int]
nums = [1, 2, 3, 4, 5]

res = map squareNum nums -- [1, 4, 9, 16, 25]
            \end{lstlisting}

            The function \texttt{squareNum} squares an integer. The list \texttt{nums} contains several integers. The \texttt{map} function is then used to apply \texttt{squareNum} to every element in \texttt{nums}, resulting in a new list \texttt{res} with each number squared.
	\section{Comparison with Other Languages}
	\section{Simple Code Examples}
	\section{Conclusion}
	\section{References}
	\section{Additional Resources}
	\section{Acknowledgments}
\end{document}
